\section{2. Gestión del proyecto con
Scrum}\label{gestiuxf3n-del-proyecto-con-scrum}

\subsection{2.1 Acuerdos de trabajo}\label{acuerdos-de-trabajo}

Los acuerdos completos están en CONTRIBUTING.md del repositorio. En
resumen:

\textbf{• Comunicación:} canal del equipo (por confirmar) para lo
urgente y GitHub (issues y Pull Requests) para todo lo técnico. Tiempo
máximo de respuesta: 24 horas en días hábiles.

\textbf{• Reuniones:} Daily breve de lunes a viernes (presencial o por
mensaje) y los eventos de Scrum al inicio y al cierre de cada sprint de
dos semanas.

\textbf{• Decisiones:} por consenso; si no se logra, decide el Product
Owner en temas de alcance y prioridad, y el Scrum Master en temas de
proceso.

\textbf{• Ramas:} main está protegida; cada integrante trabaja en
rama/\textless NOMBRE\textgreater{} y se integra solo mediante Pull
Request con al menos una revisión de otro integrante.

\textbf{• Commits:} en español, en presente y con referencia a la
historia (por ejemplo, «Agrega checklist de recepción (HU-01)»).

\textbf{• Incumplimiento:} se registra en la retrospectiva; a la segunda
falta en el mismo sprint, el Scrum Master lo informa al profesor.

\subsection{2.2 Definición de Listo (DoR) y Definición de Terminado
(DoD)}\label{definiciuxf3n-de-listo-dor-y-definiciuxf3n-de-terminado-dod}

{\def\LTcaptype{none} % do not increment counter
\begin{longtable}[]{@{}
  >{\raggedright\arraybackslash}p{(\linewidth - 2\tabcolsep) * \real{0.5001}}
  >{\raggedright\arraybackslash}p{(\linewidth - 2\tabcolsep) * \real{0.4999}}@{}}
\toprule\noalign{}
\begin{minipage}[b]{\linewidth}\raggedright
\textbf{Definición de Listo}
\end{minipage} & \begin{minipage}[b]{\linewidth}\raggedright
\textbf{Definición de Terminado}
\end{minipage} \\
\midrule\noalign{}
\endhead
\bottomrule\noalign{}
\endlastfoot
• La historia sigue el formato Como\ldots{} quiero\ldots{} para\ldots{}
y tiene ID e issue en el tablero.

• Tiene al menos 2 criterios de aceptación en formato
Dado/Cuando/Entonces.

• El equipo la estimó en puntos y cabe en un sprint (8 puntos o menos).

• Sus reglas de negocio (RN) y requisitos (RF/RNF) relacionados están
identificados.

• Sus dependencias (otras historias, datos, wireframe) están resueltas o
planificadas.

• El Product Owner la priorizó (MoSCoW). & • El código está en main
mediante un Pull Request revisado por otro integrante.

• Cumple todos sus criterios de aceptación y el Product Owner lo
verificó.

• Tiene pruebas automatizadas y la integración continua pasa sin
errores.

• Se probó solo con datos ficticios; no hay datos personales reales ni
contraseñas en el repositorio.

• La documentación afectada (SRS, modelos, README) está actualizada.

• Está desplegada en el ambiente de pruebas (desde el Sprint 4). \\
\end{longtable}
}

\subsection{2.3 Resumen de sprints}\label{resumen-de-sprints}

{\def\LTcaptype{none} % do not increment counter
\begin{longtable}[]{@{}
  >{\raggedright\arraybackslash}p{(\linewidth - 10\tabcolsep) * \real{0.0856}}
  >{\raggedright\arraybackslash}p{(\linewidth - 10\tabcolsep) * \real{0.1495}}
  >{\raggedright\arraybackslash}p{(\linewidth - 10\tabcolsep) * \real{0.3270}}
  >{\raggedright\arraybackslash}p{(\linewidth - 10\tabcolsep) * \real{0.1601}}
  >{\raggedright\arraybackslash}p{(\linewidth - 10\tabcolsep) * \real{0.1389}}
  >{\raggedright\arraybackslash}p{(\linewidth - 10\tabcolsep) * \real{0.1389}}@{}}
\toprule\noalign{}
\begin{minipage}[b]{\linewidth}\raggedright
\textbf{Sprint}
\end{minipage} & \begin{minipage}[b]{\linewidth}\raggedright
\textbf{Fechas}
\end{minipage} & \begin{minipage}[b]{\linewidth}\raggedright
\textbf{Objetivo del sprint}
\end{minipage} & \begin{minipage}[b]{\linewidth}\raggedright
\textbf{¿Se cumplió? (sí / parcial / no)}
\end{minipage} & \begin{minipage}[b]{\linewidth}\raggedright
\textbf{Puntos comprometidos}
\end{minipage} & \begin{minipage}[b]{\linewidth}\raggedright
\textbf{Puntos terminados}
\end{minipage} \\
\midrule\noalign{}
\endhead
\bottomrule\noalign{}
\endlastfoot
S0 & Por registrar & Conocer al cliente, levantar los requisitos
iniciales y preparar el repositorio y el documento base. & En curso &
--- & --- \\
S1 & & & & & \\
S2 & & & & & \\
S3 & & & & & \\
S4 & & & & & \\
S5 & & & & & \\
S6 & & & & & \\
\end{longtable}
}

\subsubsection{Evidencia de los eventos}\label{evidencia-de-los-eventos}

{\def\LTcaptype{none} % do not increment counter
\begin{longtable}[]{@{}
  >{\raggedright\arraybackslash}p{(\linewidth - 8\tabcolsep) * \real{0.0962}}
  >{\raggedright\arraybackslash}p{(\linewidth - 8\tabcolsep) * \real{0.2136}}
  >{\raggedright\arraybackslash}p{(\linewidth - 8\tabcolsep) * \real{0.2137}}
  >{\raggedright\arraybackslash}p{(\linewidth - 8\tabcolsep) * \real{0.3162}}
  >{\raggedright\arraybackslash}p{(\linewidth - 8\tabcolsep) * \real{0.1604}}@{}}
\toprule\noalign{}
\begin{minipage}[b]{\linewidth}\raggedright
\textbf{Sprint}
\end{minipage} & \begin{minipage}[b]{\linewidth}\raggedright
\textbf{Minuta de Review}
\end{minipage} & \begin{minipage}[b]{\linewidth}\raggedright
\textbf{Retrospectiva}
\end{minipage} & \begin{minipage}[b]{\linewidth}\raggedright
\textbf{Acciones de mejora acordadas}
\end{minipage} & \begin{minipage}[b]{\linewidth}\raggedright
\textbf{¿Se aplicaron?}
\end{minipage} \\
\midrule\noalign{}
\endhead
\bottomrule\noalign{}
\endlastfoot
S0 & Pendiente (al cerrar el Sprint 0) & Pendiente (al cerrar el Sprint
0) & --- & --- \\
& & & & \\
& & & & \\
\end{longtable}
}

\subsection{2.4 Registro de riesgos}\label{registro-de-riesgos}

{\def\LTcaptype{none} % do not increment counter
\begin{longtable}[]{@{}
  >{\raggedright\arraybackslash}p{(\linewidth - 10\tabcolsep) * \real{0.0856}}
  >{\raggedright\arraybackslash}p{(\linewidth - 10\tabcolsep) * \real{0.2990}}
  >{\raggedright\arraybackslash}p{(\linewidth - 10\tabcolsep) * \real{0.1282}}
  >{\raggedright\arraybackslash}p{(\linewidth - 10\tabcolsep) * \real{0.1069}}
  >{\raggedright\arraybackslash}p{(\linewidth - 10\tabcolsep) * \real{0.2735}}
  >{\raggedright\arraybackslash}p{(\linewidth - 10\tabcolsep) * \real{0.1067}}@{}}
\toprule\noalign{}
\begin{minipage}[b]{\linewidth}\raggedright
\textbf{ID}
\end{minipage} & \begin{minipage}[b]{\linewidth}\raggedright
\textbf{Riesgo}
\end{minipage} & \begin{minipage}[b]{\linewidth}\raggedright
\textbf{Probabilidad}
\end{minipage} & \begin{minipage}[b]{\linewidth}\raggedright
\textbf{Impacto}
\end{minipage} & \begin{minipage}[b]{\linewidth}\raggedright
\textbf{Respuesta y responsable}
\end{minipage} & \begin{minipage}[b]{\linewidth}\raggedright
\textbf{Estado}
\end{minipage} \\
\midrule\noalign{}
\endhead
\bottomrule\noalign{}
\endlastfoot
R-01 & Un integrante abandona el curso a mitad del semestre & Media &
Alto & Mitigar: conocimiento compartido mediante revisión cruzada de PR;
responsable: SM & Abierto \\
R-02 & La operación sin conexión y la sincronización posterior (RNF-08,
RNF-09) resultan más complejas de lo estimado & Alta & Alto & Mitigar:
prototipo técnico en el Sprint 1 y limitar lo offline a recepción y
diagnóstico; responsable: equipo de desarrollo & Abierto \\
R-03 & El cliente no está disponible para validar requisitos o asistir a
las Reviews & Media & Alto & Mitigar: agendar sesiones con anticipación
y validar por correo cuando no pueda asistir; responsable: PO &
Abierto \\
R-04 & Cambios de alcance frecuentes por parte del cliente & Media &
Medio & Mitigar: registrar cada solicitud de cambio (2.5) con análisis
de impacto; responsable: PO & Abierto \\
R-05 & El servicio gratuito de correo limita el envío de enlaces y
códigos 2FA & Media & Alto & Mitigar: probar los límites del plan
gratuito en el Sprint 1 y dejar SMS como opción (RF-02.3); responsable:
equipo de desarrollo & Abierto \\
R-06 & Las fotografías de recepción saturan el almacenamiento del
hospedaje gratuito & Media & Medio & Mitigar: comprimir las imágenes en
el dispositivo y fijar un tamaño máximo por foto; responsable: equipo de
desarrollo & Abierto \\
R-07 & Poca experiencia del equipo con FastAPI y PWA & Alta & Medio &
Mitigar: capacitación en S0--S1 y trabajo en pares en las primeras
historias; responsable: SM & Abierto \\
R-08 & Conflictos al editar en paralelo el documento .docx en ramas
distintas & Alta & Medio & Mitigar: un solo integrador por entrega y
avances por sección en docs/; responsable: SM & Abierto \\
R-09 & Se publican datos personales reales en el repositorio público &
Media & Alto & Evitar: solo datos ficticios y revisión de datos
personales en cada PR; responsable: todos & Abierto \\
\end{longtable}
}

\subsection{2.5 Solicitudes de cambio}\label{solicitudes-de-cambio}

{\def\LTcaptype{none} % do not increment counter
\begin{longtable}[]{@{}
  >{\raggedright\arraybackslash}p{(\linewidth - 10\tabcolsep) * \real{0.1282}}
  >{\raggedright\arraybackslash}p{(\linewidth - 10\tabcolsep) * \real{0.1068}}
  >{\raggedright\arraybackslash}p{(\linewidth - 10\tabcolsep) * \real{0.2137}}
  >{\raggedright\arraybackslash}p{(\linewidth - 10\tabcolsep) * \real{0.2309}}
  >{\raggedright\arraybackslash}p{(\linewidth - 10\tabcolsep) * \real{0.2137}}
  >{\raggedright\arraybackslash}p{(\linewidth - 10\tabcolsep) * \real{0.1067}}@{}}
\toprule\noalign{}
\begin{minipage}[b]{\linewidth}\raggedright
\textbf{ID / issue}
\end{minipage} & \begin{minipage}[b]{\linewidth}\raggedright
\textbf{Fecha}
\end{minipage} & \begin{minipage}[b]{\linewidth}\raggedright
\textbf{Solicitud del cliente}
\end{minipage} & \begin{minipage}[b]{\linewidth}\raggedright
\textbf{Impacto (requisitos, modelos, pruebas, esfuerzo)}
\end{minipage} & \begin{minipage}[b]{\linewidth}\raggedright
\textbf{Decisión y justificación}
\end{minipage} & \begin{minipage}[b]{\linewidth}\raggedright
\textbf{Release}
\end{minipage} \\
\midrule\noalign{}
\endhead
\bottomrule\noalign{}
\endlastfoot
--- & --- & No se han recibido solicitudes de cambio del cliente a la
fecha. & --- & --- & --- \\
\end{longtable}
}

\subsection{2.6 Declaración de uso de inteligencia
artificial}\label{declaraciuxf3n-de-uso-de-inteligencia-artificial}

{\def\LTcaptype{none} % do not increment counter
\begin{longtable}[]{@{}
  >{\raggedright\arraybackslash}p{(\linewidth - 8\tabcolsep) * \real{0.1708}}
  >{\raggedright\arraybackslash}p{(\linewidth - 8\tabcolsep) * \real{0.2137}}
  >{\raggedright\arraybackslash}p{(\linewidth - 8\tabcolsep) * \real{0.2351}}
  >{\raggedright\arraybackslash}p{(\linewidth - 8\tabcolsep) * \real{0.2307}}
  >{\raggedright\arraybackslash}p{(\linewidth - 8\tabcolsep) * \real{0.1497}}@{}}
\toprule\noalign{}
\begin{minipage}[b]{\linewidth}\raggedright
\textbf{Herramienta}
\end{minipage} & \begin{minipage}[b]{\linewidth}\raggedright
\textbf{Sección o artefacto}
\end{minipage} & \begin{minipage}[b]{\linewidth}\raggedright
\textbf{Para qué se usó}
\end{minipage} & \begin{minipage}[b]{\linewidth}\raggedright
\textbf{Cómo se verificó}
\end{minipage} & \begin{minipage}[b]{\linewidth}\raggedright
\textbf{Responsable}
\end{minipage} \\
\midrule\noalign{}
\endhead
\bottomrule\noalign{}
\endlastfoot
Claude (Anthropic), mediante Claude Code & Revisión de consistencia del
documento v0.1; borradores de 1.3, 2.1--2.4, 3.1--3.3, 3.6, 4.1 y 4.2;
correcciones de trazabilidad en 4.4--4.9; esquema de base de datos
(db/schema.sql). & Detectar inconsistencias y proponer redacción a
partir del anteproyecto, la entrevista del 19/10/2026 y el diagrama de
clases v4. & Revisión del Product Owner contra las fuentes; falta la
revisión de otro integrante antes de la entrega. & Abdías Daniel Chávez
Cruz \\
Por declarar & Secciones redactadas antes de la v0.1 y Propuesta de
Solución Tecnológica & Cada autor declara la herramienta que usó o que
no usó IA. & --- & Cada integrante \\
\end{longtable}
}
